\section*{Chapitre \ref{ch:m1:delta-one-instruments} --- Instruments delta-one}

\begin{solution}{exo:m1:delta-one-instruments:1}
Le fonds reçoit $50\,000\,000 \times 6\% = \$3\,000\,000$ et paie
$50\,000\,000 \times 4.6\% \times 91/360 = \$581\,389$~: un encaissement net de
2\,418\,611~\$.
\end{solution}

\begin{solution}{exo:m1:delta-one-instruments:2}
Spread nul~: la fair value, 6\,043.5. Cinquante points de base~: $6\,000 \times (1 +
(4.7\% - 1.3\%) \times 0.25) = 6\,051$. Chaque point d'indice sur un
contrat à terme à trois mois vaut environ 6.7~points de base de financement annuel.
\end{solution}

\begin{solution}{exo:m1:delta-one-instruments:3}
Parité $42.10 \times 2 \times 1.17 = \$98.51$~; prime $98.90/98.51 - 1 = 39$
points de base.
\end{solution}

\begin{solution}{exo:m1:delta-one-instruments:4}
Actions~: 10 de frais de négociation plus 50 d'impôt, sans financement. CFD~: 16 de frais de négociation plus 250
par an. $60 = 16 + 250\,y$ donne $y = 0.176$ année~: le CFD est moins cher pour
des durées de détention inférieures à 64 jours.
\end{solution}

\begin{solution}{exo:m1:delta-one-instruments:5}
$1.5\% \times 30\% = 45$ points de base par an~; à 15~\%, 22.5. Comme la section
871(m) assimile le paiement équivalent à un dividende d'un swap \omterm{def:m1:delta-one-instruments:deltaone}{delta-one} sur une
action américaine à un dividende de source américaine, le dealer doit prélever la \omterm{def:m1:delta-one-instruments:wht}{retenue à la source} comme sur le
dividende~: un swap qui verse 100~\% n'est pas un \omterm{def:m1:delta-one-instruments:deltaone}{delta-one}, ou bien il est mal
évalué ailleurs (un \omterm{def:m1:delta-one-instruments:swap}{spread de financement} plus large), ou repose sur une position fiscale que le détenteur
devrait souhaiter comprendre avant qu'elle ne soit contestée.
\end{solution}

\begin{solution}{exo:m1:delta-one-instruments:6}
Par an, au comptant~: les actions coûtent 30 (dividendes)~; le swap $40 + 10 = 50$~; le
contrat à terme $6 + 30 + 30 = 66$~; le fonds $7 + 30 = 37$. Sans comptant~: actions 60, swap 6,
contrat à terme 2, fonds 6. Actions contre swap~: $60 + 30y = 6 + 50y$, $y = 2.7$ années.
Fonds contre contrat à terme~: $6 + 37y = 2 + 66y$, $y = 4/29$ année, soit cinquante jours.
\end{solution}

\begin{solution}{exo:m1:delta-one-instruments:7}
Actions 80, fonds 27.75, contrat à terme 18.5, swap 18.5, CFD 86 points de base~;
annualisés 320, 111, 74, 74 et 344. Le contrat à terme et le swap sont à égalité. En pratique~: le
contrat à terme ne nécessite aucune documentation et n'implique aucune exposition à la contrepartie au-delà de la
chambre de compensation, mais n'existe que sur l'indice et en tailles fixes et
échéances~; le swap peut porter sur un panier sur mesure et reverse une plus grande part du dividende,
mais exige un contrat-cadre, une ligne de crédit et un dealer prêt à coter.
\end{solution}

\begin{solution}{exo:m1:delta-one-instruments:8}
CFD~: $16 + 3 \times 250 + 3 \times 30 = 856$ points de base. Actions~: $10 + 50 +
3 \times 30 = 150$. La page omettait le financement, facturé chaque jour sur
l'intégralité du notionnel, même à un client qui dispose des liquidités, et, en affichant «~aucune
commission~», la fourchette dans laquelle vit désormais la commission. L'affirmation est
vraie pour des durées de détention inférieures à environ deux mois.
\end{solution}

\begin{solution}{pb:m1:delta-one-instruments:1}
\textbf{1.} $5.00 \times 4 \times 1.25 = \$25.00$.
\textbf{2.} $6 + 150 + 0.05/25 = 6 + 150 + 20 = 176$ points de base.
\textbf{3.} $-(6 + 20) = -26$ points de base.
\textbf{4.} $+120$ points de base~: à l'intérieur de la bande, aucun paiement n'est dû.
\textbf{5.} La création de \omterm{def:m1:delta-one-instruments:dr}{certificats de dépôt} coûte 1.5~\% et leur annulation ne coûte pas, si bien que
l'offre réagit immédiatement aux décotes et aux primes seulement au-delà de 1.76~\%~: le
\omterm{def:m1:delta-one-instruments:dr}{certificat de dépôt} devrait se négocier avec une faible prime en moyenne, chaque fois que la demande
américaine augmente.
\textbf{6.} Acheter 400\,000 actions ordinaires~: 2\,000\,000~\pounds $=$ 2\,500\,000~\$.
Frais de dépôt 37\,500~\$, frais d'émission 5\,000~\$, négociation 1\,500~\$~: 44\,000~\$. Vendre
100\,000 ADR~: 2\,550\,000~\$. Gain 6\,000~\$, 24 points de base ($200 - 176$).
\textbf{7.} Acheter des ADR pour 2\,490\,000~\$, frais d'annulation 5\,000~\$, négociation 1\,500~\$, vendre
les actions ordinaires pour 2\,500\,000~\$~: gain 3\,500~\$, 14 points de base ($40 - 26$).
\textbf{8.} $8\% \times \sqrt{2/252} = 71$ points de base, plusieurs fois chacun des
deux gains~: la devise doit être vendue à terme au moment de la transaction. La
conversion n'est un arbitrage qu'avec trois jambes.
\textbf{9.} Les deux. Émission~: la société vend des ADR qu'elle recevra dans deux jours
et doit emprunter des ADR pour les livrer. Annulation~: elle vend des actions ordinaires
qu'elle recevra dans deux jours et doit les emprunter à Londres. Les \omterm{def:m1:financing:seclending}{frais
d'emprunt} et la disponibilité relèvent de la bande.
\textbf{10.} $25.00 \times (1 + 1.1 \times 1.5\%) = \$25.41$.
\textbf{11.} 180 points de base par rapport à la parité périmée~; 15 par rapport à l'estimation.
\textbf{12.} Il vend l'ADR et n'achète rien~: la place de Londres est fermée. Il est
vendeur à découvert de l'action pendant la nuit, couvert au mieux par des contrats à terme sur indices américains~; ce qu'il a vendu
représente 15 points de base de prime et la totalité du risque
idiosyncratique de l'action jusqu'à 08:00, heure de Londres.
\textbf{13.} Régresser le rendement de l'action ordinaire, de la clôture de Londres à la
prochaine ouverture de Londres, sur le rendement de l'ADR (ajusté du change) entre 11:30 et 16:00
heure de New York, sur quelques années, en contrôlant l'indice américain~; le coefficient
devrait être proche de 1 et le $R^2$ élevé.
\textbf{14.} Parce que chacun, dans l'enchère d'ouverture, a vu l'ADR~: le
\omterm{def:m1:auctions:imbalance}{cours indicatif} le contient déjà. La prévisibilité porte sur l'ouverture
par rapport à la clôture de la veille, et non par rapport à un prix auquel
on peut négocier.
\textbf{15.} L'ADR coûte les frais de conservation du dépositaire, quelques centimes par
certificat et par an, et (puisqu'il a été créé une seule fois) n'intègre les 1.5~\% que si
elle est celle qui le crée~; il lui évite un compte de conservation au Royaume-Uni, un règlement
en livres sterling, et la conversion de devises de chaque dividende. L'action ordinaire coûte
0.5~\% à l'achat et l'infrastructure étrangère, et elle bénéficie d'un marché plus profond.
\textbf{16.} Il ne reste plus de conversion, donc plus de bande~: une action, un
registre, la négociation à Londres devenant une cotation secondaire ou disparaissant.
\textbf{17.} Une position de relative value sans mécanisme de convergence~: elle
ne rapporte que si la restriction est levée ou si les investisseurs des deux marchés
changent d'avis. Elle devrait être dimensionnée comme un pari directionnel sur un spread.
\textbf{18.} Des dealers disposant d'un compte chez le dépositaire, une conservation dans les deux
marchés, un statut d'intermédiaire pour la fiscalité, un prêt d'actions dans les deux lignes et un
desk de change~: les \omterm{def:m1:delta-one-instruments:deltaone}{desks delta-one} des banques et quelques sociétés de trading.
\textbf{19.} \textbf{De $-26$ à $+176$ points de base.}
\textbf{20.} La même exposition dans l'enveloppe la moins chère pour le
client, évaluée à partir des différences de financement, de fiscalité et d'accès que le desk peut
supporter à moindre coût que ne le peut le client.
\end{solution}

\begin{solution}{iq:m1:delta-one-instruments:1}
Les instruments \omterm{def:m1:delta-one-instruments:deltaone}{delta-one} portent une exposition sans composante optionnelle~: actions, fonds,
contrats à terme, swaps, \omterm{def:m1:delta-one-instruments:dr}{certificats de dépôt}. Ils rapportent tous le même rendement avant frais, si bien
l'activité du desk tient aux coûts~: il cote des swaps et finance les
positions de ses clients, fait le \omterm{def:m1:what-a-trading-firm-does:market-maker}{teneur de marché} sur les fonds et les \omterm{def:m1:delta-one-instruments:dr}{certificats de dépôt}, arbitre les contrats à terme contre des paniers,
et gère les risques résiduels liés aux dividendes, à l'emprunt et au financement.
\iqlookfor{«~même exposition, coûts différents~», et au moins trois
activités concrètes.}
\end{solution}

\begin{solution}{iq:m1:delta-one-instruments:2}
Chiffrez les cinq composantes pour son horizon. Pour le fonds~: fourchette, \omterm{def:m1:the-buy-side:fees}{frais de gestion},
\omterm{def:m1:delta-one-instruments:wht}{retenue à la source} au sein du fonds. Pour le contrat à terme~: commissions infimes, quatre roll
par an, et le \omterm{def:m1:delta-one-instruments:swap}{spread de financement} implicite, qu'elle paie même si elle dispose de
liquidités. Un investisseur financé de long terme est en général mieux servi par le fonds~; un
investisseur à effet de levier ou de court terme préfère le contrat à terme~; et la réponse
dépend de la cherté du roll.
\iqlookfor{le \omterm{def:m1:delta-one-instruments:swap}{spread de financement} implicite comme frais caché du contrat à terme.}
\end{solution}

\begin{solution}{iq:m1:delta-one-instruments:3}
Les détenteurs longs de contrats à terme paient 80 points de base par an au-dessus de l'indice de référence~;
tout investisseur qui peut détenir les actions et vendre le contrat à terme le perçoit. L'opération est
cash-and-carry~: acheter le panier, vendre le contrat à terme, financer le panier. Ce qui
la freine, c'est le bilan~: la position est volumineuse, à faible rendement et consomme
du capital des dealers, surtout en fin d'année, précisément quand le \omterm{def:m1:delta-one-instruments:swap}{spread de
financement} est le plus large. Un investisseur financé peut aussi le capter passivement, en remplaçant
les contrats à terme par un fonds.
\iqlookfor{l'arbitrage \emph{et} la raison pour laquelle il est limité.}
\end{solution}

\begin{solution}{iq:m1:delta-one-instruments:4}
\omterm{def:m1:delta-one-instruments:swap}{Spread de financement}~: mon propre coût de financement de la couverture, la charge de capital et
la charge de bilan, les conditions de crédit et de marge du client, toute taxe
sur les transactions que je ne peux éviter, moins ce que me rapportent les actions en \omterm{def:m1:financing:seclending}{prêt de titres}.
Pourcentage de dividende~: mon encaissement net d'impôt sur le dividende, selon l'endroit où je
comptabilise la couverture, moins une marge. Pour une action difficile à emprunter, côté vente à
découvert, les \omterm{def:m1:financing:seclending}{frais d'emprunt}. La concurrence détermine la part de chaque avantage que je rétrocède.
\iqlookfor{bilan, position fiscale et revenus de \omterm{def:m1:financing:seclending}{prêt de titres}, les trois à la fois.}
\end{solution}

\begin{solution}{iq:m1:delta-one-instruments:5}
Seulement si la conversion est possible et moins coûteuse que 1~\%. À vérifier~: le marché d'origine
est-il ouvert (sinon la parité est périmée)~; le ratio et le taux de change retenu~;
les frais d'émission par certificat~; toute taxe de dépôt~; si le programme est ouvert
à l'émission (certains sont plafonnés ou fermés)~; l'emprunt du certificat~; le délai de
règlement-livraison et la couverture de change. Si tous les contrôles passent, émettre des certificats
adossés aux actions ordinaires.
\iqlookfor{parité périmée et cas du programme fermé.}
\end{solution}

\begin{solution}{iq:m1:delta-one-instruments:6}
Le financement devient explicite~: un spread sur le notionnel acheteur et sur le
notionnel vendeur, et non un taux d'intérêt sur la trésorerie nette. Aucune taxe à l'achat, ce qui peut
compter plus que tout le reste dans les marchés taxés. Dividendes au taux du
swap, des deux côtés. Disponibilité et frais pour la vente à découvert, tels que cotés par le
dealer, le rappel étant remplacé par le droit du dealer de résilier ou de réviser le prix.
Les réinitialisations matérialisent le P\&L en espèces chaque mois. Et la capacité dépend du bilan
du dealer, qui peut être retiré.
\iqlookfor{spreads acheteur et vendeur facturés sur le notionnel brut.}
\end{solution}
